\chapter{Introduction}
\label{cha:introduction}
\section{Background}

Recent years have seen tremendous improvements in the development of computer language models. Speech recognition and virtual assistants such as Apple's \textit{Siri}, which can be found in the company's iPhones and Smart Watches, or Amazon's \textit{Alexa}, present in products such as \textit{Echo} and \textit{Dot}, can react to human speech input with unprecedented accuracy. This trend to make  physical input devices such as keyboards, mouses or even touchscreens obsolete is likely going to gain traction and the past decade can probably be interpreted as the starting sign for further developments in this field. In the future, more and more products and services will be accessible and controllable via spoken languages only. These developments can already be observed in the automotive sector, home entertainment industries, and many more.

In order for these assistant systems to function, it is needed to translate spoken human language into structures which computers can understand.  In order to do so, these systems translate the waveforms of the spoken text into the smallest units of sound, called \textit{phonemes}. These phonemes are then combined into the corresponding words in text form and finally composed into a text representation of the spoken sentence or command. This text representation is the starting point for current machine learning algorithms to process orders such as \textit{"Hey Siri, set a ten minute timer!"} or \textit{"Hey Alexa, what events are in my calendar tomorrow?"}

Another domain that has seen a dramatic rise of practical Natural Language application are chatbots. Almost all large companies offer automated chat systems on their webpages. Most of them are in place to offer some kind of advice for selecting the right product for customers.  However, some companies are already taking the next step: Amazon for example even lets chatbots gather most of the data needed for return or service requests and only then forward customers to a human in order to finish the service.

The trend of all these developments is clear: As NLP systems are getting better and better, they will also do more and more tasks that until this day needed the skill of humans in any way. In the future, however, even more and more of these tasks are likely to be automated by language understanding systems. This development underlines the importance of NLP systems in the present and its importance on future economic processes cannot be valued enough. This is especially true because also other domains rely on NLP as a core functionality: Not only assistants and chatbots need NLP features, but also spam classification, search engines and translation systems rely on NLP in order to work properly.

\section{Objective}

This thesis is going to explore the state-of-the-art techniques which are available for extracting information from German texts. When considering the former quoted question above, it is necessary to think about what types of information are represented there. For human beings, this is quite easy to accomplish, whereas this is not the case for computers and their machine learning algorithms. In the given example above, it is quite important to recognize that the user wants to set a timer (start the timer app on a phone or smart watch) and extract the information that "ten minute" is a representation of a certain time period. Furthermore, the given information ("timer" and "ten minutes") have to be connected in a meaningful way in order to make sense and be able to be executed by the device.

\section{Thesis Structure}

This thesis demonstrates the current state of extracting information from German language texts using deep learning techniques. In order to do so, the second chapter provides the theoretical background of the applied techniques. A short introduction to the vast field of natural language processing (NLP) is followed by an explanation of \textit{word embeddings}, which is the basis of all modern text representations in the digital domain, and transforms the unstructured appearance of texts into a structured, mathematical form which can be processed by machine learning algorithms. The next section discusses the theory of \textit{transformer networks}, which represent the state of the art approach to text processing in the deep learning domain: prominent language models like \ac{BERT}, GPT2 and GPT3 are based on this network architecture and led to a mind-blowing advancement when compared to older techniques of automated text processing. The last part of the theoretical part of this paper will deal with approaches to \textit{information extraction}. Therefore, we will discuss what information can be extracted form human language: namely, we will focus on \textit{named entitiy recognition} \ac{NER}, \textit{sentiment analysis} and \textit{text summarization}. A last section of this chapter will deal with the peculiarities of the German language and what to look out for when using machine learning to extract information from German texts. 

Chapter three practically demonstrates the current state of deep learning information extraction from German language texts. In order to do so, 












